\subsection{Multiplicative representatives}

PROPOSITION 6. --- Let C be a p-ring, whose residue field k is perfect. Suppose C has finite length n (resp. infinite length). There exists a unique isomorphism \(u:W_{n}(k)\to C\) (resp. \(u:W(k)\to C\) ) which induces, by passage to quotients, the identity map of k.

Since \(W_{n}(k)\) (resp. \(W(k)\) ) is a p-ring with residue field k and of length n (resp. of infinite length) (no. 1, example 2), and since \(\varnothing\) is a p-basis of the perfect field k, prop. 6 is a special case of prop. 4 of no. 3.

THEOREM 2. --- Let A be a separated and complete local ring, k its residue field, and π the canonical homomorphism from A onto k. Suppose that k is a perfect field of characteristic p.

\begin{enumerate}
\def\labelenumi{\alph{enumi})}
\item
  There exists a unique ring homomorphism \(u: \mathbf{W}(k) \to \mathbf{A}\) such that \(\pi(u(\boldsymbol{a})) = a_0\) for \(\boldsymbol{a} = (a_n)_{n \in \mathbb{N}}\) in \(\mathbf{W}(k)\) .
\item
  The homomorphism u is continuous when W(k) is equipped with the pW(k)-adic topology, and the image of u is the unique Cohen subring of A.
\end{enumerate}

By th. 1 of no. 2, there exists a unique Cohen subring of A; denote it by C. Let u be a homomorphism from W(k) into A such that \(\pi(u(\boldsymbol{a})) = a_0\) for every \(\boldsymbol{a} = (a_n)_{n \in \mathbb{N}}\) in W(k); it is immediate that the image of u is a Cohen subring of A, hence equal to C. The existence and uniqueness of u then follow from prop. 6. The pC-adic topology of C is induced by the \(m_{A}\)-adic topology of A (no. 2, th. 1, b)), whence the continuity of u.

For a direct construction of u, see p. 70, exerc. 6.

PROPOSITION 7. --- Keep the hypotheses and notations of th. 2. There exists a unique multiplicative subset S of A such that π induces a bijection from S onto k. For an element a of A to belong to S, it is necessary and sufficient that for every n ∈ N, there exists an element \(a_{n}\) of A such that \(a = a_{n}^{p^{n}}\) . The set S is the set of elements of the form \(u(x, 0, 0, \ldots)\) .

Let us first prove the uniqueness of S. Let S be a multiplicative subset of A, such that \(\pi\) induces a bijection from S onto k. Let T be the set of elements of A which are powers \(p^{n}\) -th for every \(n \in N\) .

\begin{enumerate}
\def\labelenumi{\alph{enumi})}
\item
  We have S ⊂ T: Let \(a \in S\) and \(n \in N\) ; since the field k is perfect, there exists an element \(x_{n}\) of k such that \(x_{n}^{p^{n}} = \pi(a)\) ; since we have \(\pi(S) = k\) , there exists an element \(a_{n}\) of S such that \(x_{n} = \pi(a_{n})\) . Then one has \(\pi(a_{n}^{p^{n}}) = \pi(a)\) whence \(a_{n}^{p^{n}} = a\) since the restriction of \(\pi\) to S is injective.
\item
  The restriction of \(\pi\) to T is injective: let \(a\) and \(b\) be two elements of T such that \(\pi(a) = \pi(b)\) . Let \(n \in \mathbf{N}\) ; there exist two elements \(a_n\) and \(b_n\) of A such that \(a = a_n^{p^n}\) , \(b = b_n^{p^n}\) . Then one has \(\pi(a_n)^{p^n} = \pi(b_n)^{p^n}\) , whence \(\pi(a_n) = \pi(b_n)\) , that is, \(a_n \equiv b_n\) mod. \(\mathfrak{m}_{\mathbf{A}}\) . By lemma 1 of § 1, no. 1, we have \(a_n^{p^n} \equiv b_n^{p^n}\) mod. \(\mathfrak{m}_{\mathbf{A}}^{n+1}\) , that is, \(a \equiv b\) mod. \(\mathfrak{m}_{\mathbf{A}}^{n+1}\) . Since \(n\) is arbitrary, we have \(a = b\) .
\end{enumerate}

Properties a) and b) above, together with the formula \(\pi(S) = k\) , imply the relation \(S = T\) , whence uniqueness.

Let us now prove the existence of S. With the notations of th. 2, set \(\varphi = u \circ \tau_k\) , that is (§ 1, no. 6)

\[
\varphi (x) = u (x, 0, 0, \dots)\tag{2}
\]

for every \(x \in k\) . By prop. 4 of loc. cit., we have

\[
\varphi (1) = 1, \quad \varphi (x y) = \varphi (x) \varphi (y) \quad \text { for } \quad x, y \text {   in   } k.\tag{3}
\]

It is clear that the map \(\pi \circ \varphi\) is the identity map of \(k\) . Hence the image S of \(\varphi\) satisfies the conditions of Prop. 7.

The elements of S are often called the multiplicative (or Teichmüller) representatives of A.

Remarks. --- 1) We keep the preceding hypotheses and notation. We have

\[
\boldsymbol {a} = \sum_ {n = 0} ^ {\infty} p ^ {n} \tau_ {k} (a _ {n} ^ {p ^ {- n}}) \quad (\boldsymbol {a} = (a _ {n}) _ {n \in \mathbf {N}} \in \mathrm{W} (k))
\]

by Prop. 7 of § 1, no. 8. We therefore have

\[
u (\boldsymbol {a}) = \sum_ {n = 0} ^ {\infty} p ^ {n} \varphi (a _ {n} ^ {p ^ {- n}})\tag{4}
\]

for every \(\boldsymbol{a} = (a_n)_{n \in \mathbf{N}}\) in \(\mathrm{W}(k)\) , since \(u\) is continuous (Th. 2, b)). By formula (4), the unique Cohen subring of A consists of the elements of the form \(\sum_{n=0}^{\infty} p^n s_n\) with \(s_n \in S\) for every integer \(n \geqslant 0\) .

\begin{enumerate}
\def\labelenumi{\arabic{enumi})}
\setcounter{enumi}{1}
\tightlist
\item
  Let A be a separated and complete local ring, k its residue field, and \(\pi\) the canonical homomorphism from A onto k. One can show that there exists a multiplicative subset S of A (not unique in general) such that \(\pi\) induces a bijection from S onto k (cf. p. 72, exerc. 11).
\end{enumerate}

Examples. --- 1) Let \(k\) be a perfect field of characteristic \(p\) . The multiplicative representatives of the ring \(\mathbf{W}(k)\) are the Witt vectors \(\tau(x) = (x, 0, 0, \ldots)\) for \(x \in k\) .

\begin{enumerate}
\def\labelenumi{\arabic{enumi})}
\setcounter{enumi}{1}
\item
  Let A be an integral, separated, and complete local ring. Suppose that the residue field \(k\) of A is finite, with \(q = p^{f}\) elements, hence perfect of characteristic \(p\) . We have \(x^{q} = x\) for every \(x \in k\) , whence \(s^{q} = s\) for every multiplicative representative \(s\) . It follows that the set of multiplicative representatives consists of 0 and the \((q - 1)\)-th roots of unity in the field of fractions of A. If the field of fractions of A is locally compact, the existence of multiplicative representatives also follows from VI, § 9, no. 2, Prop. 3 (cf. also VI, § 9, exerc. 5).
\item
  More particularly, consider the case \(A = Z_{p}\) . Then the multiplicative representatives are 0 and the \((p - 1)\) -th roots of unity in the field of fractions \(Q_{p}\) of \(Z_{p}\) .
\end{enumerate}

